
\begin{ex}
	Giải hệ phương trình $\begin{cases} 3x-y=4 \\ x+2y=6 \end{cases}$.
	\loigiai{
		Ta có hệ phương trình:
		$$ \begin{cases} 3x-y=4 \quad &(1) \\ x+2y=6 \quad &(2) \end{cases} $$
		Từ phương trình (1), suy ra $y = 3x-4$.
		Thế $y = 3x-4$ vào phương trình (2):
		\begin{align*} x + 2(3x-4) &= 6 \\ x + 6x - 8 &= 6 \\ 7x &= 6 + 8 \\ 7x &= 14 \\ x &= \frac{14}{7} \\ x &= 2 \end{align*}
		Thay $x=2$ vào $y = 3x-4$:
		\begin{align*} y &= 3(2) - 4 \\ y &= 6 - 4 \\ y &= 2 \end{align*}
		Vậy nghiệm của hệ phương trình là $(x;y) = (2;2)$.
		\\ \textbf{Đáp số: $(x;y) = (2;2)$}
	}
\end{ex}

\begin{ex}
	Cho $a>0$ và $a \neq 1$. Rút gọn biểu thức $P = \left( \frac{1}{\sqrt{a}+1} - \frac{1}{a+\sqrt{a}} \right) : \frac{\sqrt{a}-1}{2\sqrt{a}+2}$.
	\loigiai{
		Điều kiện: $a>0, a \neq 1$.
		Ta có $a+\sqrt{a} = \sqrt{a}(\sqrt{a}+1)$ và $2\sqrt{a}+2 = 2(\sqrt{a}+1)$.
		\begin{align*} P &= \left( \frac{1}{\sqrt{a}+1} - \frac{1}{\sqrt{a}(\sqrt{a}+1)} \right) : \frac{\sqrt{a}-1}{2(\sqrt{a}+1)} \\ &= \left( \frac{\sqrt{a}}{\sqrt{a}(\sqrt{a}+1)} - \frac{1}{\sqrt{a}(\sqrt{a}+1)} \right) \cdot \frac{2(\sqrt{a}+1)}{\sqrt{a}-1} \\ &= \frac{\sqrt{a}-1}{\sqrt{a}(\sqrt{a}+1)} \cdot \frac{2(\sqrt{a}+1)}{\sqrt{a}-1} \end{align*}
		Vì $a>0$ và $a \neq 1$, nên $\sqrt{a}-1 \neq 0$ và $\sqrt{a}+1 \neq 0$.
		\begin{align*} P &= \frac{\cancel{(\sqrt{a}-1)}}{\sqrt{a}\cancel{(\sqrt{a}+1)}} \cdot \frac{2\cancel{(\sqrt{a}+1)}}{\cancel{(\sqrt{a}-1)}} \\ &= \frac{2}{\sqrt{a}} \end{align*}
		Vậy $P = \frac{2}{\sqrt{a}}$.
		\\ \textbf{Đáp số: $P = \frac{2}{\sqrt{a}}$}
	}
\end{ex}

\begin{ex}
	Cho phương trình $x^2-4x-1=0$, không giải phương trình, hãy tính giá trị của biểu thức $T=x_1^2+x_2^2-x_1x_2$.
	\loigiai{
		Phương trình $x^2-4x-1=0$ có $a=1, b=-4, c=-1$.
		$\Delta' = b'^2 - ac = (-2)^2 - 1(-1) = 4+1=5 > 0$.
		Phương trình có hai nghiệm phân biệt $x_1, x_2$.
		Theo hệ thức Vi-ét, ta có:
		$$ \begin{cases} x_1 + x_2 = -\frac{b}{a} = -\frac{-4}{1} = 4 \\ x_1 x_2 = \frac{c}{a} = \frac{-1}{1} = -1 \end{cases} $$
		Ta biến đổi biểu thức T:
		\begin{align*} T &= x_1^2+x_2^2-x_1x_2 \\ &= (x_1+x_2)^2 - 2x_1x_2 - x_1x_2 \\ &= (x_1+x_2)^2 - 3x_1x_2 \end{align*}
		Thay giá trị từ hệ thức Vi-ét vào biểu thức T:
		\begin{align*} T &= (4)^2 - 3(-1) \\ &= 16 + 3 \\ &= 19 \end{align*}
		Vậy $T=19$.
		\\ \textbf{Đáp số: $T=19$}
	}
\end{ex}
\begin{ex} % Câu 12 (Xác suất)
	Một chiếc hộp có chứa 5 tấm thẻ cùng loại, được đánh số lần lượt là 3; 5; 6; 7; 9. Lấy ngẫu nhiên đồng thời 2 tấm thẻ từ hộp. Tính xác suất của biến cố A: "Tích các số ghi trên 2 tấm thẻ chia hết cho 3".
	\loigiai{
		Tập hợp các kết quả có thể xảy ra khi lấy ngẫu nhiên đồng thời 2 tấm thẻ từ hộp là:
		(3;5); (3;6); (3;7); (3;9); (5;6); (5;7); (5;9); (6;7); (6;9); (7;9).
		Số phần tử của không gian mẫu là 10.
		
		Tập hợp các kết quả thuận lợi cho biến cố A (tích các số ghi trên 2 tấm thẻ chia hết cho 3) là:
		(3;5); (3;6); (3;7); (3;9); (5;6); (5;9); (6;7); (6;9); (7;9).
		Số phần tử của biến cố A là 9.
		
		Xác suất của biến cố A là:
		\begin{align*}
			P(A) = \frac{9}{10}
		\end{align*}
	}
\end{ex}

\begin{ex} % Câu 12 (Bài toán lập hệ phương trình)
	Một mảnh vườn hình chữ nhật có chiều dài hơn chiều rộng 15m. Nếu giảm chiều dài 2m và tăng chiều rộng 3m thì diện tích tăng thêm 44m$^2$. Tính diện tích mảnh vườn đó.
	\loigiai{
		Gọi chiều dài ban đầu của mảnh vườn là $x$ (m) và chiều rộng ban đầu là $y$ (m).
		Điều kiện: $x > y > 0$, $x > 2$.
		
		Theo đề bài, chiều dài hơn chiều rộng 15m, ta có phương trình:
		$x - y = 15 \quad (1)$
		
		Diện tích ban đầu của mảnh vườn là $S_1 = xy$ (m$^2$).
		
		Nếu giảm chiều dài 2m, chiều dài mới là $x-2$ (m).
		Nếu tăng chiều rộng 3m, chiều rộng mới là $y+3$ (m).
		Diện tích mới của mảnh vườn là $S_2 = (x-2)(y+3)$ (m$^2$).
		
		Theo đề bài, diện tích tăng thêm 44m$^2$, nên ta có phương trình:
		$S_2 - S_1 = 44$
		\begin{align*}
			(x-2)(y+3) - xy &= 44 \\
			xy + 3x - 2y - 6 - xy &= 44 \\
			3x - 2y - 6 &= 44 \\
			3x - 2y &= 50 \quad (2)
		\end{align*}
		Từ (1) và (2), ta có hệ phương trình:
		$$ \begin{cases} x - y = 15 \\ 3x - 2y = 50 \end{cases} $$
		Từ (1) suy ra $x = y + 15$. Thế vào (2):
		\begin{align*}
			3(y+15) - 2y &= 50 \\
			3y + 45 - 2y &= 50 \\
			y + 45 &= 50 \\
			y &= 5
		\end{align*}
		Giá trị $y=5$ thỏa mãn điều kiện $y>0$.
		Thay $y=5$ vào $x = y+15$:
		$x = 5 + 15 = 20$.
		Giá trị $x=20$ thỏa mãn điều kiện $x>y$ (vì $20>5$) và $x>2$.
		
		Vậy, chiều dài ban đầu của mảnh vườn là 20m và chiều rộng ban đầu là 5m.
		Diện tích ban đầu của mảnh vườn là:
		\begin{align*}
			S_1 = x \times y = 20 \times 5 = 100 \quad (\text{m}^2)
		\end{align*}
		Vậy diện tích mảnh vườn đó là 100 m$^2$.
	}
\end{ex}
\begin{ex} % Câu 13
	Cho đường tròn $(O;R)$, đường kính $AB$. Gọi $D$ là trung điểm của đoạn $OB$. Từ $D$ kẻ dây $PQ$ của đường tròn $(O)$ vuông góc với đường kính $AB$. Lấy $M$ là một điểm bất kỳ trên cung nhỏ $AP$, dây $MB$ cắt dây $PQ$ tại $I$.
	\begin{enumerate}
		\item Chứng minh 4 điểm $D, I, M, A$ cùng thuộc một đường tròn.
		\item Chứng minh $QB^2 = BI \cdot BM$.
		\item Gọi $C$ là một điểm nằm trên dây $MQ$ sao cho $MP = MC$. Xác định vị trí $M$ trên cung nhỏ $AP$ để tổng $S = MA + MP$ đạt giá trị lớn nhất.
	\end{enumerate}
\end{ex}

